% =============================================================
% Resultados
% Apartado individual, no es un Capitulo
% =============================================================
\chapter*{Resultados}
\label{cap:resultados}
\phantomsection
\addcontentsline{toc}{chapter}{Resultados}

A la fecha de redacci\'on de este informe, los resultados obtenidos son
parciales, dado que el proyecto se encuentra en curso (ver
Secci\'on~\ref{sec:metodologia_intro}). Se resumen a continuaci\'on los
resultados alcanzados hasta el momento:

\begin{itemize}
  \item El circuito detector de picos dise\~nado (Cap\'itulo~\ref{cap:diseno_deteccion_picos})
        fue simulado y validado experimentalmente, confirmando su
        comportamiento lineal para tensiones de entrada superiores a
        0.8~V y su correcta respuesta ante pulsos de 100~ns
        \cite{informeEneFeb2026,informeAbril2026}.
  \item Las oscilaciones observadas en el prototipo inicial, atribuidas
        a capacidades par\'asitas, fueron resueltas mediante el
        rediseño de la placa con t\'ecnicas de ruteo para se\~nales de
        alta velocidad (Secci\'on~\ref{sec:mitigacion_parasitos})
        \cite{informeMarzo2026,informeAbril2026}.
  \item Se verific\'o la correcta comunicaci\'on entre la placa de
        control y el l\'aser seeder mediante el enlace RS232/MAX3232
        \cite{informeAbril2026}.
  \item Se implement\'o un primer borrador del software HMI, con
        comunicaci\'on funcional en sentido placa-PC
        \cite{informeAbril2026}.
  \item Avances parciales del proyecto fueron presentados en el ILRC32
        \cite{galvagno2026ilrc}.
\end{itemize}

PLACEHOLDER: completar este apartado con los resultados de la
validaci\'on integral del sistema (firmware completo, comunicaci\'on
bidireccional del HMI y mediciones sobre el canal HSR real), en
relaci\'on con los objetivos espec\'ificos planteados en la
Secci\'on~\ref{sec:objetivos_especificos}, una vez finalizada la etapa
de desarrollo restante.
